\section{Fusion-Sufficient Moment Exchange}
\label{sec:method}

\subsection{Label-Wise Moment Projection}

At time $k$, source $j$ maintains an LMB posterior
\begin{equation}
\begin{split}
    \pi_j &= \{(r_j^\ell,p_j^\ell):\ell\in\mathcal L_j\},\\
    p_j^\ell(x)&=\sum_{m=1}^{M_j^\ell}w_{jm}^\ell
    \mathcal N(x;m_{jm}^\ell,P_{jm}^\ell).
\end{split}
    \label{eq:lmb-gm}
\end{equation}
For normalized $w_{jm}^\ell$, define the label-wise projection $\mathcal P$ by
\begin{equation}
 \mathcal P(\pi_j)=\{(r_j^\ell,\mathcal N(\mu_j^\ell,\Sigma_j^\ell))\}_\ell,
 \quad \mu_j^\ell=\sum_m w_{jm}^\ell m_{jm}^\ell,
 \label{eq:projection}
\end{equation}
\begin{equation}
 \Sigma_j^\ell=\sum_m w_{jm}^\ell\!\left[P_{jm}^\ell+
 (m_{jm}^\ell-\mu_j^\ell)(m_{jm}^\ell-\mu_j^\ell)^\top\right].
 \label{eq:covariance}
\end{equation}
This keeps the label and existence probability and replaces the mixture by the
Gaussian with identical first two moments. The projection is idempotent:
$\mathcal P(\mathcal P(\pi))=\mathcal P(\pi)$.

\subsection{The Projected Gaussian KLA-LMB Receiver}

Receiver $i$ first applies $\mathcal P$ to every delivered source, aligns the
union of labels, and then fuses one Gaussian per label. Let
$\mathcal S_i^\ell$ contain self and delivered neighbors carrying label
$\ell$. Let $\alpha_{ij}$ and $\beta_{ij}$ be the spatial and existence
weights. Missing-label renormalization gives
\begin{equation}
 a_{ij}^\ell=\frac{\alpha_{ij}}{\sum_{q\in\mathcal S_i^\ell}\alpha_{iq}},
 \qquad
 b_{ij}^\ell=\frac{\beta_{ij}}{\sum_{q\in\mathcal S_i^\ell}\beta_{iq}}.
 \label{eq:active-weights}
\end{equation}
The Gaussian canonical parameters are
\begin{equation}
\begin{split}
 K_i^\ell&=\sum_{j\in\mathcal S_i^\ell}a_{ij}^\ell(\Sigma_j^\ell)^{-1},
 \quad h_i^\ell=\sum_j a_{ij}^\ell(\Sigma_j^\ell)^{-1}\mu_j^\ell,\\
 \bar\Sigma_i^\ell&=(K_i^\ell)^{-1},\qquad
 \bar\mu_i^\ell=\bar\Sigma_i^\ell h_i^\ell.
\end{split}
 \label{eq:gaussian-fusion}
\end{equation}
Define the Gaussian overlap by
\begin{equation}
\begin{split}
 \eta_i^\ell&=\int \varphi_i^\ell(x)\,dx,\\
 \varphi_i^\ell(x)&=\prod_{j\in\mathcal S_i^\ell}
 \mathcal N(x;\mu_j^\ell,\Sigma_j^\ell)^{a_{ij}^\ell}.
\end{split}
 \label{eq:gaussian-overlap}
\end{equation}
The Bernoulli terms are
\begin{equation}
 q_1^\ell=\eta_i^\ell\prod_j(r_j^\ell)^{b_{ij}^\ell},\quad
 q_0^\ell=\prod_j(1-r_j^\ell)^{b_{ij}^\ell},\quad
 \bar r_i^\ell=\frac{q_1^\ell}{q_0^\ell+q_1^\ell}.
 \label{eq:existence-fusion}
\end{equation}
where all products in Eq.~\eqref{eq:existence-fusion} are over
$j\in\mathcal S_i^\ell$.
The implementation clamps each $r_j^\ell$ to $[10^{-9},1-10^{-9}]$ and
lower-bounds $q_0^\ell+q_1^\ell$ by machine epsilon; both paths use the same
guards.
Denote Eqs.~\eqref{eq:gaussian-fusion}--\eqref{eq:existence-fusion} by
$\mathcal G_\omega$, where $\omega=(\alpha,\beta)$, and the receiver by
$\mathcal F_\omega=\mathcal G_\omega\circ\mathcal P$.

\textbf{Proposition 1 (fusion-output equivalence).}
For each source, assume identical active-label presence after selection and
pruning; identical message-type-independent spatial and existence weights;
identical delivery masks; common projection and covariance regularization;
disabled covariance inflation; and a common lossless codec. Then
\begin{equation}
 \mathcal F_\omega(\pi_1,\ldots,\pi_n)=
 \mathcal F_\omega(\mathcal P\pi_1,\ldots,\mathcal P\pi_n).
 \label{eq:commutation}
\end{equation}
\emph{Proof.} Projection preserves each source-label's $r$, $\mu$, and
$\Sigma$ and does not change its presence in $\mathcal S_i^\ell$. Hence the
renormalized $a_{ij}^\ell,b_{ij}^\ell$, canonical $K_i^\ell,h_i^\ell$,
overlap $\eta_i^\ell$, and Bernoulli terms $q_0^\ell,q_1^\ell$ are unchanged.
The common codec and covariance regularizer therefore produce the same
$\bar r_i^\ell,\bar\mu_i^\ell,\bar\Sigma_i^\ell$. Equivalently,
$\mathcal G_\omega\circ\mathcal P\circ\mathcal P=
\mathcal G_\omega\circ\mathcal P$ by idempotence. $\square$

We call $\mathcal P(\pi)$ \emph{fusion-sufficient} only with respect to this
$\mathcal F_\omega$: messages in the same equivalence class yield the same
fusion output. Equation~\eqref{eq:commutation} does not imply equal posterior
densities. It need not hold if a receiver preserves mixture modes, adapts
weights from higher-order structure, uses different projection regularization,
or reuses a compressed result in a multi-round algorithm.

\subsection{Typed Message and Byte Accounting}

Both arms use the same versioned, little-endian binary codec: a 24-byte message
header, a 20-byte label header, and binary64 component fields for weight, mean,
and the upper covariance triangle. For state dimension $d$, a message with
$M_\ell$ components per label has
\begin{equation}
 B=24+\sum_{\ell}\left[20+8M_\ell
 \left(1+d+\frac{d(d+1)}{2}\right)\right]\ \text{bytes}.
 \label{eq:wire-bytes}
\end{equation}
The full message retains every $M_\ell$ in Eq.~\eqref{eq:lmb-gm}; the proposed
message sets $M_\ell=1$ using Eqs.~\eqref{eq:projection}--\eqref{eq:covariance}.
Attempted bytes are counted after encoding and before the drop draw; delivered
bytes count messages decoded after successful delivery. No other tracker
variable, scheduling decision, or fusion parameter changes.
